\chapter{Introduction}

\section{Biological Context for Computer Scientists}
\Glspl{protein} are the basic building blocks of a cell and its DNA \gls{genome} encodes the information needed to build them. Every cell of a multicellular organism carries essentially the same DNA, yet cells develop into distinct types. The difference cannot come from the DNA itself. Between the reading of the genome and the building of proteins, therefore, lies one or more layers of regulation that determine which proteins a cell produces, and those proteins in turn give the cell its characteristics. Studing these layers is central to explaining how cells function and what shapes their behaviour.

Gene expression proceeds in three stages. A \gls{gene} is first copied into a long molecule of \gls{rna}, called \gls{unsplicedTranscript}. The second stage, \gls{splicing}, removes the non-coding blocks and joins the coding ones at the resulting \glspl{spliceJunction}, yielding a shorter spliced \gls{transcript}. The final, spliced version is \gls{mrna}, which the ribosome translates\defn{Translation}{The process by which the ribosome reads an mRNA molecule and assembles the corresponding chain of amino acids into a protein.} into protein in the third and final phase.

The shortening process is not a one-to-one mapping. To expand the decoding capacity of a gene, cells can produce several different spliced versions of the same \gls{unsplicedTranscript}; these versions are called \glspl{isoform}. Because different isoforms can translate into proteins with entirely different, or even conflicting, roles, a cell can fundamentally alter its behaviour simply by switching which variant it expresses---the MET$\Delta$14 isoform, for instance, acts as a switch that drives invasive growth in cancer \cite{cerqua2022met}. The core goal of \gls{dtu} analysis is therefore to measure changes in the relative expression proportions of specific transcript isoforms from the same gene across different biological conditions, developmental stages, or cell types, in order to expose regulatory events that gene-level measurement renders invisible.

\section{A Brief History of Single-Cell Sequencing}
For most of its history, \gls{rnaSeq} was performed on bulk samples. RNA is extracted from a population of cells, pooled into a single reaction, and sequenced as one library\defn{Sequencing library}{The set of DNA molecules loaded into the sequencer for one run. Here it holds the sample's RNA, copied into DNA and tagged so that the instrument can read each fragment.}, so the resulting profile describes the sample as a whole rather than any cell within it. Sequencing of this kind, termed \gls{bulkRnaSeq}, is statistically comfortable: pooling thousands to millions of cells yields deep coverage and stable abundance estimates for every gene. The averaging, however, is also its principal limitation. A gene's measured abundance is the mean over all cells present, so cellular heterogeneity\defn{Cellular heterogeneity}{The coexistence of distinct cell types and states within a single organism.} disappears before the data ever reach an analyst. A rare cell type constituting two percent of a tissue contributes two percent of that tissue's \glspl{read} for each of its markers\defn{Marker genes}{A gene or transcript whose expression specifically identifies a cell type, state, or condition.} and is easily buried beneath the dominant population; two subpopulations that shift in opposite directions---one gaining a transcript while the other loses it---cancel out and appear unchanged.

\gls{scRnaSeq} replaced that aggregate view with per-cell profiling, running the same measurement independently on thousands of cells. The strategy rests on compartmentalisation. Each cell is physically isolated, its polyadenylated RNA captured and reverse-transcribed separately, and every resulting molecule marked with two labels. A \gls{cellBarcode} identifies the cell of origin, so reads can be regrouped into per-cell profiles after sequencing; a \gls{umi} tags the individual molecule, so \gls{pcr} duplicates are collapsed rather than miscounted as expression. A single pooled sequencing run is thereby converted into thousands of independent expression measurements, each far shallower than a bulk sample.

The first \gls{scRnaSeq}, produced from 2009 onwards, came from plate-based protocols, which consume one cell per reaction \cite{tang2009mrna}. The chemistry of that era anchored its reads at one end of the molecule, and base accuracy declines along a read as it lengthens.\defn{Base}{One of the four nucleotides that form a DNA or RNA strand: adenine (A), cytosine (C), guanine (G), and thymine (T). A read is a sequence of consecutive bases, and the confidence with which each base is identified declines along a read as it lengthens.} Reliable information therefore concentrated near that end, a limitation that deeper sequencing could not remedy---data of this kind exhibit \gls{tagBias} \cite{hashimshony2012cel}.

In 2012, the invention of \gls{fullLength} removed that bias \cite{ramskold2012full, picelli2013smart}. A whole molecule could now be copied and read from end to end, preserving the internal structure that distinguishes one transcript variant from another. While this was a substantial gain in coverage, the technology was costly and the plate format kept a study at hundreds to a few thousand cells.

A new protocol, droplet microfluidics, suddenly revolutionised the scale in 2015 \cite{zheng2017massively}: each droplet encapsulates one cell together with a bead\defn{Bead}{A small gel sphere loaded with the barcoded primers. One bead is co-encapsulated with each cell, supplying every primer its droplet needs.} carrying barcoded primers\defn{Barcoded primers}{Primers that carry a short identifying sequence, so that every molecule captured in a given droplet can be traced back to its cell after sequencing.}, so tens of thousands of cells are profiled in one run at a small fraction of the cost per cell. The gain came at a price. Droplet capture is anchored at the transcript tail; this forces the use of the old chemistry, and the tag bias returned.

\Gls{dropletProtocols} dominate single-cell work today, since they profile far more cells per experiment at a far lower cost per cell than plate-based chemistry allows \cite{conte2024opportunities}. What tag bias removes is therefore the central limitation of the format rather than an incidental one, and recovering it is an active line of work. Long-read platforms read whole transcripts at single-cell resolution, and short-read datasets are mined for whatever isoform signal their terminal fragments still carry \cite{joglekar2023words}.

\section{The Computational Challenge}
Turning droplet data into transcript-level results means working around two limitations that the sequencing chemistry itself imposes, one on the quantity of data and one on its content. The first is sparsity. \Glspl{countMatrix} from droplet experiments are mostly zeros, with zero entries routinely forming 70\% to 90\% of all values. Some of those zeros are biological, recording a transcript that the cell never produced, while others are technical, recording a molecule that was present but never captured, and both look identical in the matrix.

The second limitation is positional. Only the terminal \gls{threePrimeEnd} of each molecule is read, so where several \glspl{isoform} share that end they look identical in the data and their proportions cannot be measured directly. Deeper sequencing reduces the first limitation and cannot touch the second: extra reads fill in more of the matrix, but no quantity of terminal fragments recovers the interior of a transcript. A method that works on droplet data must therefore extract signal from a lightly observed count matrix while accepting that most of each molecule was never seen.

\subsection*{A Note on What Remains Recoverable}

Gene-level information survives the \gls{tagBias} constraint. A gene registers when a read is confidently attributed to it, and in most cases a read that cannot be traced to a single transcript still maps to isoforms of the same gene. A gene matrix can therefore be built from the same reads, and it remains informative for \gls{deg} analysis. A consequential property is that this matrix is denser than the transcript matrix.

\section{Current Approaches and Limitations}\label{sec:intro-approaches}
Transcript-level regulation has not gone unstudied; the question is whether existing tools survive the sparsity and positional bias that droplet data impose. Three families of methods address it, and each stops short of detecting \gls{dtu} across a full droplet dataset.

\textbf{Bulk algorithms.} Established \gls{bulkRnaSeq} tools such as rMATS \cite{shen2014rmats}, SUPPA2 \cite{trincado2018suppa}, DRIMSeq \cite{nowicka2016drimseq} and DEXSeq \cite{anders2012dexseq} model continuous isoform proportions from deep, junction-spanning coverage. Applied to sparse droplet matrices, those assumptions break down and the resulting estimates are unstable.

\textbf{Aggregation to pseudo-bulk.} One route around the coverage problem is \gls{pseudoBulk}, which sums the cells of a cluster or condition into a single profile to recover the depth that bulk tools expect \cite{gilis2021satuRn, tekath2021dturtle}. The cost is the cell-to-cell variation the assay was built to measure: a switch confined to a small subpopulation is diluted by the cells around it.

\textbf{Targeted \gls{threePrime} tools.} Sierra \cite{patrick2020sierra} analyses droplet data directly, but only for a narrow class of events at transcript termini. It offers no genome-wide framework for multi-\gls{isoformSwitch} from raw zero-heavy counts.

What is missing is a framework that works on sparse droplet matrices at single-cell resolution, without collapsing cells together and without full-length reads.

\section{Thesis Objective}

The question this thesis pursues originates in COTAN (Co-expression Tables Analysis) \cite{galfre2021cotan}, a framework for gene-level \glspl{countMatrix} that focuses its statistical model on the zeros of the matrix rather than on its positive counts. COTAN asks whether two genes are detected together less often than chance would predict, using presence/absence contingency tables, so no artificial imputation\defn{Imputation}{Replacing zero counts with values predicted from neighbouring cells, at the risk of inventing spurious correlations.} or variance-stabilising transformation\defn{Variance-stabilising transformation}{A transformation applied to counts, such as a logarithm, so that their variance no longer depends on their mean. It alters the scale of the data, turning an undetected molecule into a small positive value rather than a zero.} is needed, and it performs well on the sparse matrices typical of single-cell data. The present thesis asks whether the same statistics hold when applied one level lower, to \gls{transcript} matrices whose counts are sparser still: whether they retain enough \gls{isoform} information to detect \gls{dtu} without fixing the \gls{tagBias} problem. The claim is deliberately modest --- isoforms that share their terminal region remain indistinguishable --- and the aim is to map where sparsity-robust statistics recover differential transcript usage anyway.

The work proceeded from raw reads to a reported result in four stages, each producing a concrete artefact: a count-matrix pipeline, an analysis library, a reproducible driver layer, and a set of reported results.

\textbf{From reads to transcript matrices.} A Nextflow pipeline was built that turns raw droplet reads into a filtered transcript-level count matrix---the artefact the whole study rests on, since DTU analysis requires transcript resolution that gene-level matrices do not carry.

\textbf{An analysis library.} The analysis was consolidated into a single R package, \texttt{cotanisoform}, which extends COTAN's gene-level framework to transcript resolution. It keeps COTAN's statistics---co-expression, p-values, GDI, differential expression---but adds the steps a transcript-level analysis needs. It is where COTAN's gene-framed statistics become a transcript-level detector.

\textbf{The analysis drivers.} A configuration-driven layer reproduces a full per-dataset analysis from a single declarative file, joining the pipeline and the library so the whole study runs end to end reproducibly.

\textbf{From candidates to reported results.} The workflow was applied to two datasets---a homogeneous human cell-line mixture and complex mouse cortical tissue---and the candidates are cross-referenced against published isoform biology and gene-level clustering baselines. The analysis was carried out first-hand and not reviewed by a domain expert; the findings are a reproducible first report rather than established biology, and independent verification is expected before they are built upon.

